\documentclass[a4paper,12pt]{report}
\usepackage[utf8]{inputenc}
\usepackage[T1]{fontenc}
\usepackage[french]{babel}
\usepackage{geometry}
\usepackage{graphicx}
\usepackage{booktabs}
\usepackage{tikz}
\usetikzlibrary{shapes,arrows,positioning,fit,calc}
\usepackage{algorithm}
\usepackage{algpseudocode}
\usepackage{amsmath}
\usepackage{amssymb}
\usepackage{caption}
\usepackage{subcaption}
\usepackage{xcolor}
\usepackage{hyperref}

\geometry{margin=2.5cm}

\title{Chapitre 2 : Modélisation CAD réaliste et intégration URDF}
\author{Rapport de stage ingénieur en robotique}
\date{Août 2026}

\begin{document}

\maketitle

\tableofcontents

\chapter{Modélisation CAD réaliste et intégration URDF}

\section{Introduction}

\subsection{Motivation du passage au modèle CAD réaliste}

Le Chapitre 1 a présenté un prototype initial basé sur des primitives géométriques simples (boîtes et cylindres) pour valider les concepts fondamentaux de la cinématique tracteur-remorque et de la navigation ROS2. Bien que cette approche ait permis de démontrer la viabilité du système, plusieurs limitations ont motivé le passage à un modèle CAD réaliste :

\begin{itemize}
    \item \textbf{Précision physique insuffisante} : Les primitives géométriques ne représentent pas fidèlement la distribution de masse réelle du robot, ce qui affecte le comportement dynamique en simulation.
    \item \textbf{Matrices d'inertie approximatives} : Les formes simples ne permettent pas d'estimer correctement les moments d'inertie, cruciaux pour la stabilité du contrôle et la simulation physique réaliste.
    \item \textbf{Placement des capteurs imprécis} : Les positions et orientations des capteurs (LIDAR, IMU, GPS) par rapport au centre de masse ne correspondent pas à la réalité physique.
    \item \textbf{Validation limitée} : Un modèle simplifié ne permet pas de tester les scénarios réels d'interaction avec l'environnement (collision, franchissement d'obstacles).
    \item \textbf{Transition vers prototype réel} : L'objectif final étant le déploiement sur un robot physique, un modèle CAD réaliste facilite le transfert vers la fabrication.
\end{itemize}

Ce chapitre décrit donc le processus complet de modélisation CAD réaliste, depuis la sélection des modèles 3D jusqu'à l'intégration fonctionnelle dans l'écosystème ROS2/Gazebo via URDF/xacro.

\section{Recherche et sélection des modèles CAD}

\subsection{Critères de sélection}

La recherche de modèles CAD appropriés s'est appuyée sur les critères suivants :

\begin{table}[htbp]
\centering
\caption{Critères de sélection des modèles CAD}
\begin{tabular}{@{}llp{8cm}@{}}
\toprule
\textbf{Critère} & \textbf{Poids} & \textbf{Description} \\
\midrule
Format d'export & Critique & Compatibilité avec ROS2/Gazebo (STL, STEP préférentiel) \\
Complexité géométrique & Élevé & Représentativité des formes réelles du robot \\
Distribution de masse & Critique & Permet le calcul précis des matrices d'inertie \\
Échelle et dimensions & Élevé & Correspondance avec les spécifications du prototype \\
Licence d'utilisation & Moyen & Modèles libres de droits pour usage académique/industriel \\
Facilité de modification & Moyen & Possibilité d'ajuster les géométries si nécessaire \\
\bottomrule
\end{tabular}
\end{table}

\subsection{Composants identifiés}

Les modèles CAD suivants ont été sélectionnés pour constituer l'assemblage robot + remorque :

\begin{table}[htbp]
\centering
\caption{Composants CAD sélectionnés}
\begin{tabular}{@{}lllp{6cm}@{}}
\toprule
\textbf{Composant} & \textbf{Type} & \textbf{Source} & \textbf{Caractéristiques} \\
\midrule
Châssis tracteur & Assembly & SolidWorks & Structure principale avec compartiments électroniques \\
Roues motrices & Part & SolidWorks & Pneus avec jante, rayon 0.085m \\
Roues castor & Part & SolidWorks & Roulettes avant orientables \\
Châssis remorque & Assembly & SolidWorks & Plateforme de chargement avec essieu \\
Roues remorque & Part & SolidWorks & Roues fixes arrière remorque \\
Support capteurs & Part & SolidWorks & Structure montée sur châssis principal \\
LIDAR YDLidar & Part & Fournisseur & Modèle X4 avec spécifications officielles \\
IMU MPU-9250 & Part & Fournisseur & Boîtier avec orientation repère \\
GPS u-blox & Part & Fournisseur & Antenne avec boîtier d'étanchéité \\
\bottomrule
\end{tabular}
\end{table}

\section{Extraction et modification CAD/STL}

\subsection{Pipeline SolidWorks vers URDF}

Le processus de conversion des modèles CAO vers un format compatible ROS2 suit le pipeline illustré Figure~\ref{fig:cad_pipeline}.

\begin{figure}[htbp]
\centering
\begin{tikzpicture}[
    node distance=1.5cm,
    auto,
    block/.style={rectangle, draw, fill=blue!10, text width=2.5cm, text centered, rounded corners, minimum height=1.2cm},
    arrow/.style={->, >=stealth, thick}
]

\node [block] (solidworks) {SolidWorks};
\node [block, right=of solidworks] (export) {Export STEP/STL};
\node [block, right=of export] (meshlab) {Simplification maillage};
\node [block, right=of meshlab] (stl) {Fichiers STL};
\node [block, below=of stl] (urdf) {Intégration URDF};
\node [block, left=of urdf] (xacro) {Macros xacro};
\node [block, left=of xacro] (gazebo) {Simulation Gazebo};

\draw [arrow] (solidworks) -- (export);
\draw [arrow] (export) -- (meshlab);
\draw [arrow] (meshlab) -- (stl);
\draw [arrow] (stl) -- (urdf);
\draw [arrow] (urdf) -- (xacro);
\draw [arrow] (xacro) -- (gazebo);

\node [above=0.2cm of solidworks, font=\footnotesize] {Modèles CAO};
\node [above=0.2cm of export, font=\footnotesize] {Format d'échange};
\node [above=0.2cm of meshlab, font=\footnotesize] {Optimisation};
\node [above=0.2cm of stl, font=\footnotesize] {Maillages};
\node [below=0.2cm of urdf, font=\footnotesize] {Définition robot};
\node [below=0.2cm of xacro, font=\footnotesize] {Paramétrage};
\node [below=0.2cm of gazebo, font=\footnotesize] {Validation};

\end{tikzpicture}
\caption{Pipeline de conversion CAD vers URDF}
\label{fig:cad_pipeline}
\end{figure}

\subsection{Export depuis SolidWorks}

L'export depuis SolidWorks a suivi les étapes suivantes :

\begin{enumerate}
    \item \textbf{Préparation de l'assemblage} : Suppression des composants non pertinents (vis, écrous, détails internes) pour réduire la complexité.
    \item \textbf{Définition des repères} : Création de repères (coordinate systems) pour chaque sous-assemblage correspondant aux liens URDF futurs.
    \item \textbf{Export STL} : Export de chaque composant individuel en format binaire STL pour compatibilité Gazebo.
    \item \textbf{Vérification des échelles} : Confirmation que les dimensions exportées correspondent aux spécifications métriques (mm $\to$ m conversion).
\end{enumerate}

\subsection{Simplification des maillages}

Les fichiers STL bruts issus de SolidWorks présentent souvent une densité de triangles excessive pour la simulation en temps réel. L'outil \textit{MeshLab} a été utilisé pour :

\begin{itemize}
    \item Réduire le nombre de faces tout préservant la forme globale (réduction typique de 70\%)
    \item Réparer les erreurs de maillage (trous, faces non-manifold)
    \item Unifier les normales pour un rendu cohérent
    \item Optimiser la taille des fichiers pour un chargement rapide
\end{itemize}

\subsection{Organisation des fichiers maillages}

Les fichiers STL finaux sont organisés dans la structure suivante :

\begin{table}[htbp]
\centering
\caption{Organisation des fichiers maillages}
\begin{tabular}{@{}llp{7cm}@{}}
\toprule
\textbf{Chemin} & \textbf{Fichier} & \textbf{Description} \\
\midrule
meshes/visual/ & chassis.stl & Maillage haute résolution pour rendu visuel \\
meshes/collision/ & chassis\_collision.stl & Maillage simplifié pour physique \\
meshes/visual/ & wheel.stl & Roue avec détails de pneu \\
meshes/collision/ & wheel\_collision.stl & Roue cylindre simplifié \\
meshes/visual/ & trailer.stl & Châssis remorque complet \\
meshes/sensors/ & lidar.stl & Boîtier LIDAR \\
meshes/sensors/ & imu.stl & Boîtier IMU \\
meshes/sensors/ & gps.stl & Antenne GPS \\
\bottomrule
\end{tabular}
\end{table}

\section{Construction URDF/xacro}

\subsection{Architecture xacro}

L'utilisation du langage \textit{xacro} (XML Macro) permet une définition modulaire et paramétrable du modèle robot. L'architecture des fichiers est illustrée Figure~\ref{fig:xacro_architecture}.

\begin{figure}[htbp]
\centering
\begin{tikzpicture}[
    node distance=1.2cm,
    auto,
    file/.style={rectangle, draw, fill=green!10, text width=2.8cm, text centered, rounded corners, minimum height=0.8cm},
    main/.style={rectangle, draw, fill=red!10, text width=3cm, text centered, rounded corners, minimum height=1cm},
    arrow/.style={->, >=stealth, thick}
]

\node [main] (main) {robot\_remorque.urdf.xacro};
\node [file, above left=of main] (chassis) {chassis.xacro};
\node [file, above=of main] (wheels) {wheels.xacro};
\node [file, above right=of main] (trailer) {trailer.xacro};
\node [file, below left=of main] (sensors) {sensors\_gps\_imu.xacro};
\node [file, below=of main] (control) {tractor\_ros2\_control.xacro};
\node [file, below right=of main] (materials) {materials.xacro};

\draw [arrow] (main) -- (chassis);
\draw [arrow] (main) -- (wheels);
\draw [arrow] (main) -- (trailer);
\draw [arrow] (main) -- (sensors);
\draw [arrow] (main) -- (control);
\draw [arrow] (main) -- (materials);

\end{tikzpicture}
\caption{Architecture des fichiers xacro}
\label{fig:xacro_architecture}
\end{figure}

\subsection{Tableau des liens et joints}

Le modèle URDF définit l'arbre cinématique complet du système tracteur-remorque. Le Tableau~\ref{tab:links_joints} répertorie les liens principaux et leurs connexions.

\begin{table}[htbp]
\centering
\caption{Liens et joints du modèle URDF}
\label{tab:links_joints}
\begin{tabular}{@{}lllllp{4cm}@{}}
\toprule
\textbf{Link} & \textbf{Joint} & \textbf{Type} & \textbf{Parent} & \textbf{Child} & \textbf{Description} \\
\midrule
base\_footprint & base\_footprint\_joint & fixed & - & base\_link & Frame racine au sol \\
base\_link & - & - & base\_footprint & - & Châssis tracteur \\
rear\_left & rear\_left\_joint & continuous & base\_link & rear\_left & Roue motrice gauche \\
rear\_right & rear\_right\_joint & continuous & base\_link & rear\_right & Roue motrice droite \\
front\_left & front\_left\_joint & continuous & base\_link & front\_left & Caster avant gauche \\
front\_right & front\_right\_joint & continuous & base\_link & front\_right & Caster avant droit \\
support & support\_joint & fixed & base\_link & support & Support capteurs \\
lidar\_link & lidar\_joint & fixed & support & lidar\_link & LIDAR YDLidar \\
camera\_link & camera\_joint & fixed & support & camera\_link & Caméra \\
imu\_link & imu\_joint & fixed & support & imu\_link & IMU \\
gps\_link & gps\_joint & fixed & support & gps\_link & GPS \\
trailer\_base\_link & hitch\_joint & revolute & base\_link & trailer\_base\_link & Attelage remorque \\
trailer\_wheel\_left & trailer\_wheel\_left\_joint & continuous & trailer\_base\_link & trailer\_wheel\_left & Roue remorque gauche \\
trailer\_wheel\_right & trailer\_wheel\_right\_joint & continuous & trailer\_base\_link & trailer\_wheel\_right & Roue remorque droite \\
\bottomrule
\end{tabular}
\end{table}

\subsection{Maillages visuels vs collision}

Chaque lien URDF possède deux représentations géométriques distinctes :

\begin{table}[htbp]
\centering
\caption{Comparaison maillages visuels et collision}
\begin{tabular}{@{}llp{5cm}p{5cm}@{}}
\toprule
\textbf{Aspect} & \textbf{Maillage visuel} & \textbf{Maillage collision} & \textbf{Justification} \\
\midrule
Résolution & Haute (détails) & Basse (simplifié) & Performance physique \\
Fichiers STL & meshes/visual/*.stl & meshes/collision/*.stl & Séparation claire \\
Complexité & Formes réalistes & Primitives simples & Stabilité simulation \\
Couleurs & Textures/Matériaux & Couleur unie & Visualisation vs physique \\
Calcul & Rendu graphique & Détection collision & Optimisation \\
\bottomrule
\end{tabular}
\end{table}

\subsection{Matrices d'inertie}

Les matrices d'inertie sont calculées à partir de la géométrie CAD en utilisant les propriétés de masse de SolidWorks. Pour chaque lien, les paramètres suivants sont définis :

\begin{table}[htbp]
\centering
\caption{Paramètres d'inertie des liens principaux}
\begin{tabular}{@{}llccp{4cm}@{}}
\toprule
\textbf{Link} & \textbf{Masse (kg)} & \textbf{IXX (kg·m²)} & \textbf{IYY (kg·m²)} & \textbf{Centre de masse (m)} \\
\midrule
base\_link & 12.5 & 0.45 & 0.52 & (0.0, 0.0, 0.085) \\
rear\_left & 1.2 & 0.008 & 0.012 & (0.0, 0.0, 0.0) \\
rear\_right & 1.2 & 0.008 & 0.012 & (0.0, 0.0, 0.0) \\
trailer\_base\_link & 8.3 & 0.32 & 0.38 & (0.25, 0.0, 0.0) \\
trailer\_wheel\_left & 0.8 & 0.005 & 0.008 & (0.0, 0.0, 0.0) \\
trailer\_wheel\_right & 0.8 & 0.005 & 0.008 & (0.0, 0.0, 0.0) \\
\bottomrule
\end{tabular}
\end{table}

La matrice d'inertie complète pour un lien est définie comme :

\begin{equation}
\mathbf{I} = \begin{bmatrix}
I_{xx} & I_{xy} & I_{xz} \\
I_{yx} & I_{yy} & I_{yz} \\
I_{zx} & I_{zy} & I_{zz}
\end{bmatrix}
\end{equation}

Pour les liens avec symétrie géométrique (roues), les produits d'inertie hors-diagonale sont négligés ($I_{xy} = I_{xz} = I_{yz} = 0$).

\section{Intégration des capteurs réels}

\subsection{Placement TF des capteurs}

Les capteurs sont positionnés dans l'arbre TF avec des offsets précis par rapport au châssis principal. Le Tableau~\ref{tab:sensor_tf} détaille ces positions.

\begin{table}[htbp]
\centering
\caption{Placement TF des capteurs}
\label{tab:sensor_tf}
\begin{tabular}{@{}lllp{5cm}@{}}
\toprule
\textbf{Capteur} & \textbf{Frame} & \textbf{Parent} & \textbf{Position (x,y,z) en m} \\
\midrule
LIDAR & lidar\_link & support & (-0.02, 0.04, 0.20) \\
IMU imu\_link & support & (0.0, 0.0, 0.05) \\
GPS & gps\_link & support & (0.2, 0.0, 0.1) \\
Caméra & camera\_link & support & (0.08, 0.0, 0.14) \\
\bottomrule
\end{tabular}
\end{table}

\subsection{Plugins Gazebo}

Chaque capteur est associé à un plugin Gazebo spécifique pour simuler son comportement physique et publier les données ROS2 correspondantes.

\begin{table}[htbp]
\centering
\caption{Plugins Gazebo pour les capteurs}
\begin{tabular}{@{}lllp{5cm}@{}}
\toprule
\textbf{Capteur} & \textbf{Type Gazebo} & \textbf{Plugin} & \textbf{Topic ROS2} \\
\midrule
LIDAR & gpu\_lidar & libignition-sensors-gpu-lidar & /scan \\
IMU & imu & libignition-sensors-imu & /imu \\
GPS & navsat & libignition-sensors-navsat & /gps/fix \\
Caméra & camera & libignition-rendering-ogre2 & /image\_raw \\
\bottomrule
\end{tabular}
\end{table}

\subsection{Spécifications des capteurs}

\begin{table}[htbp]
\centering
\caption{Spécifications techniques des capteurs simulés}
\begin{tabular}{@{}llp{6cm}@{}}
\toprule
\textbf{Capteur} & \textbf{Paramètre} & \textbf{Valeur} \\
\midrule
LIDAR & Portée min/max & 0.15 / 12.0 m \\
& Nombre d'échantillons & 240 \\
& Champ de vision & ±120° \\
& Fréquence & 10 Hz \\
\midrule
IMU & Fréquence & 50 Hz \\
& Bruit gyro & 0.01 rad/s \\
& Bruit accéléro & 0.1 m/s² \\
\midrule
GPS & Fréquence & 10 Hz \\
& Écart-type position & 0.5 m \\
& Écart-type vitesse & 0.1 m/s \\
\midrule
Caméra & Résolution & 640×480 \\
& Champ devision & 90° \\
& Fréquence & 30 Hz \\
\bottomrule
\end{tabular}
\end{table}

\section{Stack technologique de simulation}

\subsection{Outils et versions}

La simulation s'appuie sur une stack technologique cohérente et documentée, résumée dans le Tableau~\ref{tab:tech_stack}.

\begin{table}[htbp]
\centering
\caption{Stack technologique de simulation}
\label{tab:tech_stack}
\begin{tabular}{@{}llp{6cm}@{}}
\toprule
\textbf{Catégorie} & \textbf{Outil} & \textbf{Rôle/Version} \\
\midrule
CAD source & SolidWorks & Modélisation 3D, export STEP/STL \\
Format d'échange & STEP/STL & Format standard d'interopérabilité CAO \\
Simplification maillage & MeshLab & Optimisation des fichiers STL \\
ROS2 & Humble & Middleware robotique (2024 LTS) \\
Simulateur & Gazebo Ignition Fortress & Simulation physique temps réel \\
Bridging & ros\_gz\_bridge & Communication ROS2/Gazebo \\
Contrôle & gz\_ros2\_control & Interface ros2_control pour Gazebo \\
Langage URDF & xacro & Macros XML pour modèles paramétrables \\
Navigation & Nav2 & Stack de navigation ROS2 \\
Localisation & robot\_localization & Fusion multi-capteurs (EKF) \\
\bottomrule
\end{tabular}
\end{table}

\subsection{Comparaison Gazebo Classic vs Ignition}

Le passage de Gazebo Classic (Chapitre 1) à Gazebo Ignition (Chapitre 2) introduit des changements significatifs dans l'architecture de simulation.

\begin{table}[htbp]
\centering
\caption{Comparaison Gazebo Classic vs Ignition}
\begin{tabular}{@{}llp{4cm}p{4cm}@{}}
\toprule
\textbf{Aspect} & \textbf{Gazebo Classic} & \textbf{Gazebo Ignition} & \textbf{Impact} \\
\midrule
Architecture & Monolithique & Modulaire (librairies) & Flexibilité accrue \\
Plugins & C++ spécifiques Gazebo & Ignition Sensors/Rendering & Portabilité améliorée \\
Bridging ROS2 & ros\_gazebo\_bridge & ros\_gz\_bridge & Stabilité temporelle \\
Horloge /clock & Intégrée & Séparée, nécessite bridge & Synchronisation explicite \\
Format monde & SDF 1.6+ & SDF 1.9+ & Nouvelles fonctionnalités \\
Physique & ODE & Ignition Physics & Précision améliorée \\
Contrôle & gazebo\_ros\_control & gz\_ros2\_control & Interface standardisée \\
\bottomrule
\end{tabular}
\end{table}

\subsection{Challenges de migration}

La migration vers Gazebo Ignition a présenté plusieurs challenges techniques :

\begin{itemize}
    \item \textbf{Synchronisation de l'horloge} : Le topic /clock doit être explicitement bridgé entre Gazebo et ROS2, contrairement à Gazebo Classic où cette intégration était native.
    \item \textbf{Plugins capteurs} : Les plugins Gazebo Classic ne sont pas compatibles ; il faut utiliser les plugins Ignition Sensors avec une syntaxe différente.
    \item \textbf{Format SDF} : Certaines balises SDF ont changé de nom ou de structure, nécessitant une adaptation des fichiers monde.
    \item \textbf{Documentation} : La documentation d'Ignition étant moins mature, le débogage a nécessité plus d'investigation expérimentale.
\end{itemize}

\section{Architecture logicielle}

\subsection{Arborescence xacro}

L'organisation modulaire des fichiers xacro permet une maintenance aisée et une réutilisation des composants. La Figure~\ref{fig:xacro_tree} illustre cette arborescence.

\begin{figure}[htbp]
\centering
\begin{tikzpicture}[
    level distance=1.2cm,
    sibling distance=2.5cm,
    edge from parent/.style={draw, -latex},
    every node/.style={draw, rectangle, rounded corners, align=center, fill=blue!5}
]

\node {robot\_remorque.urdf.xacro}
    child { node {chassis.xacro} }
    child { node {wheels.xacro} }
    child { node {trailer.xacro} }
    child { node {sensors\_gps\_imu.xacro} }
    child { node {tractor\_ros2\_control.xacro} }
    child { node {materials.xacro} };

\end{tikzpicture}
\caption{Arborescence des fichiers xacro}
\label{fig:xacro_tree}
\end{figure}

\subsection{Arbre TF du système}

L'arbre des transformations (TF) du système complet tracteur-remorque est illustré Figure~\ref{fig:tf_tree}.

\begin{figure}[htbp]
\centering
\begin{tikzpicture}[
    level distance=1cm,
    sibling distance=1.5cm,
    edge from parent/.style={draw, -latex},
    every node/.style={draw, rectangle, rounded corners, align=center, font=\footnotesize, fill=green!5}
]

\node {map}
    child { node {odom}
        child { node {base\_footprint}
            child { node {base\_link}
                child { node {support}
                    child { node {camera\_link} }
                    child { node {lidar\_link} }
                    child { node {imu\_link} }
                    child { node {gps\_link} }
                }
                child { node {rear\_left} }
                child { node {rear\_right} }
                child { node {front\_left} }
                child { node {front\_right} }
            }
            child { node {trailer\_base\_link}
                child { node {trailer\_wheel\_left} }
                child { node {trailer\_wheel\_right} }
            }
        }
    };

\end{tikzpicture}
\caption{Arbre TF du système tracteur-remorque}
\label{fig:tf_tree}
\end{figure}

\subsection{Organisation des packages ROS2}

Le projet est structuré en deux packages principaux avec des responsabilités clairement définies.

\begin{table}[htbp]
\centering
\caption{Organisation des packages ROS2}
\begin{tabular}{@{}llp{6cm}@{}}
\toprule
\textbf{Package} & \textbf{Type} & \textbf{Contenu} \\
\midrule
diff\_robot & ament\_cmake & Modèles URDF/xacro, launch files, configurations, contrôleurs \\
trailer\_kinematics & ament\_python & Nœuds de cinématique remorque (EKF beta, footprint) \\
\bottomrule
\end{tabular}
\end{table}

\subsection{Structure des dossiers}

\begin{table}[htbp]
\centering
\caption{Structure des dossiers du package diff\_robot}
\begin{tabular}{@{}llp{6cm}@{}}
\toprule
\textbf{Dossier} & \textbf{Sous-dossiers} & \textbf{Description} \\
\midrule
config/ & *.yaml & Fichiers de configuration (Nav2, contrôleurs, EKF) \\
launch/ & *.py & Fichiers de lancement ROS2 \\
urdf/ & *.xacro, *.urdf, meshes/ & Modèles robot et maillages STL \\
world/ & *.world & Mondes Gazebo SDF \\
rviz/ & *.rviz & Configurations RViz2 \\
map/ & *.yaml, *.pgm & Cartes statiques pour navigation \\
\bottomrule
\end{tabular}
\end{table}

\section{Algorithmes et méthodes de calcul implémentés}

\subsection{Calcul des matrices d'inertie depuis la géométrie CAD}

Les matrices d'inertie sont estimées à partir des propriétés de masse exportées depuis SolidWorks. L'algorithme suivant formalise ce processus.

\begin{algorithm}[htbp]
\caption{Estimation des matrices d'inertie depuis CAD}
\label{alg:inertia_estimation}
\begin{algorithmic}[1]
\State \textbf{Entrée:} Modèle CAO avec propriétés de masse pour chaque lien
\State \textbf{Sortie:} Matrices d'inertie URDF pour chaque lien
\For{chaque lien $l$ dans l'assemblage}
    \State Extraire masse $m_l$ depuis SolidWorks
    \State Extraire centre de masse $\mathbf{c}_l = (c_x, c_y, c_z)$
    \State Extraire moments d'inertie $I_{xx}, I_{yy}, I_{zz}$ au centre de masse
    \State Calculer produits d'inertie $I_{xy}, I_{xz}, I_{yz}$ (si asymétrie)
    \State Définir matrice d'inertie $\mathbf{I}_l$ dans repère lien
    \State Ajuster $\mathbf{I}_l$ si centre de masse $\neq$ origine lien (théorème de Huygens)
\EndFor
\State \textbf{Retourner} ensemble des matrices $\{\mathbf{I}_l\}$
\end{algorithmic}
\end{algorithm}

\subsection{Placement des frames capteurs (calibration extrinsèque)}

Le positionnement précis des capteurs nécessite une calibration extrinsèque définissant les transformations entre le repère du capteur et le repère du robot.

\begin{algorithm}[htbp]
\caption{Algorithme de placement des frames capteurs}
\label{alg:sensor_placement}
\begin{algorithmic}[1]
\State \textbf{Entrée:} Spécifications géométriques robot, positions capteurs cibles
\State \textbf{Sortie:} Transformations TF pour chaque capteur
\For{chaque capteur $s$}
    \State Définir frame parent $P_s$ (généralement base\_link ou support)
    \State Mesurer/extraire position $\mathbf{p}_s = (p_x, p_y, p_z)$ dans repère $P_s$
    \State Définir orientation $\mathbf{R}_s$ (quaternion ou rpy)
    \State Créer joint $j_s$ de type \textit{fixed} entre $P_s$ et frame\_$s$
    \State Spécifier transformation $\mathbf{T}_s = (\mathbf{p}_s, \mathbf{R}_s)$ dans URDF
    \State Vérifier cohérence avec spécifications constructeur
\EndFor
\State \textbf{Retourner} ensemble des transformations $\{\mathbf{T}_s\}$
\end{algorithmic}
\end{algorithm}

\subsection{Suivi de trajectoire GPS (waypoint following)}

Pour tester le châssis réaliste en navigation, un algorithme de suivi de waypoints basé sur GPS a été implémenté.

\begin{algorithm}[htbp]
\caption{Algorithme de suivi de waypoints GPS}
\label{alg:waypoint_following}
\begin{algorithmic}[1]
\State \textbf{Entrée:} Liste de waypoints $\{(lat_i, lon_i)\}$, position GPS courante
\State \textbf{Sortie:} Commande vitesse $(v, \omega)$
\State Convertir waypoints actuel et courant en coordonnées cartésiennes locales
\State Calculer vecteur cible $\mathbf{d} = \mathbf{p}_{target} - \mathbf{p}_{current}$
\State Calculer distance au waypoint $d = \|\mathbf{d}\|$
\If{$d < d_{threshold}$}
    \State Passer au waypoint suivant
    \State Si tous les waypoints visités, arrêter
\EndIf
\State Calculer angle de cap désiré $\theta_{des} = \arctan2(d_y, d_x)$
\State Calculer erreur angulaire $\Delta\theta = \theta_{des} - \theta_{current}$
\State Appliquer contrôleur proportionnel :
\State $v = K_v \cdot \min(d, d_{max})$
\State $\omega = K_\omega \cdot \Delta\theta$
\State \textbf{Retourner} $(v, \omega)$
\end{algorithmic}
\end{algorithm}

\section{Validation du modèle dans Gazebo}

\subsection{Stabilité physique}

La validation du modèle CAD réaliste a révélé plusieurs aspects critiques de la stabilité physique :

\begin{itemize}
    \item \textbf{Centrede masse} : Le centre de masse calculé depuis CAD a permis une stabilité améliorée par rapport au modèle primitives, réduisant les basculements lors des accélérations.
    \item \textbf{Matrices d'inertie} : Les moments d'inertie réalistes ont produit un comportement de rotation plus fidèle, particulièrement visible lors des virages serrés.
    \item \textbf{Contact roue-sol} : La géométrie détaillée des roues a amélioré le modèle de contact, réduisant les glissements intempestifs.
\end{itemize}

\subsection{Cohérence TF}

L'arbre TF a été validé en utilisant l'outil \texttt{tf2\_tools} :

\begin{itemize}
    \item Vérification de la connectivité complète de l'arbre
    \item Confirmation des fréquences de publication (50 Hz pour les joints, 10 Hz pour les capteurs)
    \item Validation des transformations statiques vs dynamiques
    \item Test de latence des transformations (typiquement < 5 ms)
\end{itemize}

\subsection{Difficultés résolues}

Plusieurs difficultés ont été rencontrées et résolues lors de l'intégration :

\begin{table}[htbp]
\centering
\caption{Difficultés rencontrées et solutions}
\begin{tabular}{@{}llp{6cm}@{}}
\toprule
\textbf{Problème} & \textbf{Cause} & \textbf{Solution} \\
\midrule
Robot instable au spawn & Centre de masse incorrect & Recalcul depuis CAD, ajustement offset \\
Maillages trop lourds & Trop de triangles STL & Simplification MeshLab (réduction 70\%) \\
Capteurs non publiés & Plugins Gazebo incompatibles & Migration vers plugins Ignition Sensors \\
TF tree incomplet & Joints manquants & Ajout joints fixed pour capteurs \\
Sauts horloge /clock & Bridge non configuré & Ajout explicite bridge clock dans launch \\
Collision imprécise & Maillages visuels utilisés & Création maillages collision simplifiés \\
\bottomrule
\end{tabular}
\end{table}

\subsection{Métriques de validation}

Les métriques suivantes ont été utilisées pour valider le modèle :

\begin{table}[htbp]
\centering
\caption{Métriques de validation du modèle}
\begin{tabular}{@{}llp{6cm}@{}}
\toprule
\textbf{Métrique} & \textbf{Valeur cible} & \textbf{Résultat} \\
\midrule
Temps de chargement monde & < 10 s & 8.3 s \\
Fréquence simulation & ≥ 1000 Hz & 1200 Hz (stable) \\
Latence TF & < 10 ms & 3-5 ms \\
Stabilité robot (test accélération) & Pas de basculement & Stable jusqu'à 1.5 m/s² \\
Précision odométrie & < 5\% erreur & 3.2\% erreur sur 10 m \\
\bottomrule
\end{tabular}
\end{table}

\section{Conclusion}

Ce chapitre a présenté le processus complet de modélisation CAD réaliste et d'intégration URDF pour le système tracteur-remorque. Les principaux accomplissements sont :

\begin{itemize}
    \item \textbf{Modèle CAD réaliste} : Conversion réussie des modèles SolidWorks en maillages STL optimisés pour Gazebo.
    \item \textbf{URDF/xacro fonctionnel} : Architecture modulaire permettant une maintenance aisée et une réutilisation des composants.
    \item \textbf{Intégration capteurs} : Placement précis des capteurs (LIDAR, IMU, GPS) avec plugins Gazebo appropriés.
    \item \textbf{Stack technologique cohérente} : Migration réussie de Gazebo Classic vers Ignition avec adaptation des plugins et bridges.
    \item \textbf{Validation physique} : Stabilité améliorée du modèle grâce aux matrices d'inertie réalistes.
\end{itemize}

Le modèle CAD réaliste constitue désormais une base solide pour les développements futurs, notamment :

\begin{itemize}
    \item La fusion de données multi-capteurs pour la localisation robuste (Chapitre 4)
    \item L'expérimentation de scénarios de navigation complexes en environnement réaliste
    \item La préparation au transfert vers un prototype physique
\end{itemize}

Le chapitre suivant abordera la cinématique avancée de la remorque et les algorithmes de contrôle spécifiques aux systèmes tracteur-remorque.

\end{document}
